%==============================================================================
\section{How the results were checked}\label{app:verify}
%==============================================================================

This appendix says how the results of this report were checked, so that a reader who
wants to repeat the check can.

\subsection{How the program works}\label{ssec:verify-principle}

It shares no code with the solvers, the campaign harness, or any script written during
the research, so where the two agree, that is a second opinion rather than the same
calculation run twice. It also avoids using the results it is checking. The loading rule
of Proposition~\ref{prop:ktns} is implemented separately and compared against an exact
dynamic program over magazine states rather than assumed optimal, and the optimum $\Zst$
comes from a dynamic program over subsets of jobs, cross-checked against brute-force
enumeration on the smallest instances, which uses none of the methods of
Section~\ref{sec:exact}. On small instances it enumerates rather than samples; where a
family is too large to enumerate, the sample and its generation parameters are given and
the statement is labelled an Observation rather than a Proposition.

\subsection{What is covered}\label{ssec:verify-covered}

\begin{itemize}[topsep=4pt,itemsep=2pt]
  \item \textbf{The loading oracle.} KTNS against the exact dynamic program, on
  $349{,}872$ instance-and-sequence pairs, including every permutation of every instance
  in a random bank. No disagreement. The same comparison was run against the
  implementation used by the solvers, with the same result.
  \item \textbf{The convention identity} of Proposition~\ref{prop:conv}, on random
  instances.
  \item \textbf{Both lower bounds} and Corollary~\ref{cor:lb}, on every instance in the
  random census.
  \item \textbf{Every named instance} in Sections~\ref{sec:prelim} and \ref{sec:gap}:
  the $k$-ring for $k=3,\dots,9$ in both readings; the sliding-window ring; the
  witness $W$ of Proposition~\ref{prop:refute}; the pair $I_0$, $I_1$ of
  Proposition~\ref{prop:noclutter}; the $\Kst=4$ instance; the smallest sub-optimal
  instance of Observation~\ref{obs:smallest}; the copies of
  Observation~\ref{obs:unbounded-true}; the star and non-matroid instances of
  Propositions~\ref{prop:chromK} and \ref{prop:nonmatroid}.
  \item \textbf{Every bound in Section~\ref{ssec:ratio} and \ref{ssec:law}}, checked
  for violations on a random census of $200$ instances: Theorem~\ref{thm:uncond},
  Corollaries~\ref{cor:zerogap}, \ref{cor:smallZ} and \ref{cor:z3},
  Lemma~\ref{lem:transcap}, and Propositions~\ref{prop:k3} and \ref{prop:genk}. No
  violation.
  \item \textbf{The closed form of Proposition~\ref{prop:hk3}}, on instances with
  $\Kst=3$ drawn at random until $45$ were obtained. No disagreement.
  \item \textbf{The covering relaxation} of Proposition~\ref{prop:oddring}, by solving
  the fractional cover as a linear program for $k=3,\dots,8$. This is what identified
  that the statement requires $k\ge4$.
  \item \textbf{The setup-cost collapse} of Theorem~\ref{thm:collapse}, by enumerating
  all feasible groupings and all orderings of an instance and minimising the augmented
  objective directly. This is also what identified that the threshold must be read in
  the walk value.
  \item \textbf{The PCF relaxation} of Proposition~\ref{prop:pcf}, by building the
  linear program explicitly over all configurations and solving it, with and without
  the counting rows.
\end{itemize}

\subsection{What the verification found}\label{ssec:verify-found}

Two findings changed the content of this report.

The first is the distinction of Section~\ref{ssec:frames}. The $k$-ring family is
treated in the literature as the canonical example of a positive gap for the two-phase
heuristic, with the gap unbounded on copies of it. The verification found that every
$k$-ring has \emph{zero} gap for the method as specified, and that the published values
are walk values. That is why
Section~\ref{ssec:frames} exists, why Table~\ref{tab:rings} has two gap columns, and
why Proposition~\ref{prop:unbounded} is now stated for the walk value with
Observation~\ref{obs:unbounded-true} supplying a different family for the other
reading.

The second is Proposition~\ref{prop:oddring}, which had been stated for all $k$ and
fails at $k=3$, where the whole job set forms a single feasible group.

Both were errors of statement rather than of proof: the arguments were sound for the
objects they were actually about.

A third finding concerns the campaign analysis rather than the structural results. Two
properties of the raw result files can silently corrupt any analysis of them; both are
in the code that reads the results rather than in any solver, and
Appendix~\ref{app:repro} states them.

\subsection{Where the checking stops}\label{ssec:verify-limits}

The checking is exhaustive only on small instances. On larger ones it compares methods
against each other, which catches disagreement but not an error they share. And
reproducing a number confirms the number rather than the argument given for it, so the
proofs were re-derived by hand as well.

\subsection{Running it}\label{ssec:verify-running}

The program is \texttt{verification/verify\_report\_independent.py} in the repository.
It requires Python~3 and SciPy, needs no commercial solver, and completes in roughly
twenty-five minutes. It prints one line per claim, naming the section of this report the
claim comes from, the value that section states, and the value the computation
produces, and exits with a non-zero status if any of them disagree.

The campaign numbers of Section~\ref{sec:comp} are produced by a separate analysis over
the raw per-instance result files, described in Appendix~\ref{app:repro}. No number in
Section~\ref{sec:comp} is carried over from an earlier campaign.
